\section{\textbf{Related Work}}

\subsection{\textbf{General Intelligence and Cognitive Architectures}}

Research on general machine intelligence predates contemporary deep learning. Turing shifted the question of whether a machine could think toward an operational test based on observable interaction \cite{turing1950computing}. Later definitions sought to characterize intelligence independently of a particular embodiment or task. Legg and Hutter formalized universal intelligence in terms of an agent's expected performance across computable environments \cite{legg2007universal}, while Goertzel described AGI as the construction of systems possessing broad, general-purpose intelligence rather than competence restricted to a narrow domain \cite{goertzel2014agi}. Chollet subsequently argued that task-specific skill is insufficient as a measure of intelligence and emphasized skill-acquisition efficiency, priors, experience, and generalization to novel problems \cite{chollet2019measure}. These formulations differ in their assumptions and evaluation methods, but collectively move the discussion away from performance on a fixed task and toward adaptive competence across tasks and environments.

Architectural approaches to this goal historically emphasized explicit cognitive organization. Soar proposed an implemented architecture for general intelligent behavior based on problem spaces, goals, production rules, learning, and persistent knowledge \cite{laird1987soar}. ACT-R similarly integrated perceptual--motor processing, goals, declarative memory, production rules, and subsymbolic selection mechanisms within a unified theory of cognition \cite{anderson2004integrated}. These systems established an important principle retained by AGI-Former: general behavior is unlikely to arise from an isolated prediction function without mechanisms that connect perception, memory, internal state, decision, and action. Their integration, however, is primarily constructed through symbolic or hybrid cognitive modules. AGI-Former investigates a different implementation path in which multimodal representations and causal thought--action generation are connected through an end-to-end trainable deep neural architecture.

\subsection{\textbf{Generalist and Multimodal Sequence Models}}

Transformer-based generalist models replace some hand-designed cognitive interfaces with shared sequence modeling. Gato demonstrated that a single network with a common token interface could perform language, image captioning, simulated control, and real-robot tasks \cite{reed2022generalist}. Perceiver IO generalized attention-based processing across structured input and output spaces through a latent representation \cite{jaegle2022perceiver}. Flamingo and PaLM-E showed complementary methods for connecting visual or continuous sensor representations to large language models \cite{alayrac2022flamingo,driess2023palme}. As discussed in Section~I, these works establish that heterogeneous modalities, tasks, and embodiments can share substantial computation. For the present comparison, their most important distinction is architectural: Gato is organized as a generalist sequence policy, Perceiver IO uses a latent input--output interface, and Flamingo and PaLM-E remain centered on language-model computation conditioned by non-textual representations.

More direct perception-to-action systems have narrowed the distance between multimodal modeling and agency. VIMA encodes interleaved visual and textual prompts and autoregressively generates robot-control actions, demonstrating multimodal prompting and zero-shot generalization across manipulation tasks \cite{jiang2023vima}. RT-2 casts robotic actions as tokens and co-trains vision--language data with robot trajectories, allowing an end-to-end vision-language-action model to transfer semantic knowledge into robotic control \cite{zitkovich2023rt2}. These systems are particularly relevant because action is part of the learned output rather than an external interpretation of generated language. Their scope is nevertheless centered on embodied manipulation with visual and linguistic conditioning; they do not attempt the proposed organization of several modality-native encoders around a persistent, causal thought--action state.

\subsection{\textbf{Unified Multimodal Encoder--Decoder Models}}

Unified-IO 2 is the closest existing architectural reference point for AGI-Former. It converts images, text, audio, actions, bounding boxes, and other supported inputs and outputs into tokens in a shared semantic space and processes them with a single encoder--decoder Transformer \cite{lu2024unifiedio2}. Joint training across a large multimodal mixture enables one model to perform understanding, generation, and robotic-manipulation tasks. This result provides strong evidence that encoder--decoder Transformers can support broad multimodal input and output spaces without requiring a separate complete model for every task.

The similarity is substantial but does not make the architectures equivalent. Unified-IO 2 primarily formulates heterogeneous tasks as autoregressive multimodal sequence generation within a shared token space. AGI-Former instead begins from an agent-centered decomposition. Modality-specific encoders preserve the native organization of image, video, audio, language, and other sensory streams; prior thought and action tokens form the causal state of the decoder; and cross-attention determines how the current state selects evidence from each modality. The proposed disjoint organization retains separate modality-conditioned outputs when actions have incompatible execution semantics. The joint organization first performs action-conditioned selection within each modality and then fuses the selected representations to generate a coordinated action stream. Thus, the intended contribution is not multimodal token unification alone, but an explicit recurrent relationship among perception, thought, and action.

This distinction also produces different experimental questions. Unified multimodal models are commonly evaluated through aggregate performance over many understanding and generation benchmarks. AGI-Former additionally requires evaluation of whether an evolving thought--action state improves cross-modal selection, whether late fusion improves coordinated decisions, whether modality-specific encoders retain information lost through universal tokenization, and whether the same learned organization transfers across environments, embodiments, and action spaces. These questions allow the architectural claims to be tested through controlled ablations rather than inferred from the breadth of supported tasks.

\subsection{\textbf{Model Orchestration and Modular Agents}}

A second line of work achieves broad capability by coordinating existing models rather than jointly training their internal representations. HuggingGPT uses a language model for task planning, expert-model selection, execution, and response synthesis across several modalities \cite{shen2023hugginggpt}. This approach is closely related to the disjoint AGI-Former baseline, in which modality specialists communicate through language with a supervising language model. Modular orchestration offers practical advantages: individual components can be replaced, independently improved, or selected according to cost and capability. It also creates explicit interfaces at which behavior can be inspected or constrained.

The same modularity introduces a different learning boundary. Routing decisions, representation conversion, tool invocation, and recovery from component failure are implemented partly outside the participating neural models. Consequently, the final action objective does not necessarily train every perceptual and reasoning component through one differentiable path. AGI-Former does not treat modular orchestration as an obsolete design; it uses it as an implementable comparison point and as a possible stage in architectural development. The empirical question is whether deeper joint training produces measurable gains in cross-modal grounding, coordination, transfer, efficiency, or robustness that justify its reduced component independence.

\subsection{\textbf{Positioning of AGI-Former}}

Prior work has therefore established five necessary foundations: formal accounts of general intelligence, integrated cognitive architectures, generalist sequence policies, unified multimodal encoder--decoder models, and perception-to-action learning. No single distinction among these categories is sufficient to establish AGI, and AGI-Former does not claim otherwise. Its specific architectural hypothesis is that integrated multimodal agency can be studied through modality-native encoding, a persistent causal thought--action state, and direct action-conditioned interaction between that state and sensory representations.

Relative to symbolic cognitive architectures, AGI-Former replaces predetermined representational interfaces with learned attention operations. Relative to decoder-centered generalist models, it restores an explicit encoder--decoder relationship between perception and action. Relative to language-conditioned multimodal models, it does not require every modality to become text before joint reasoning. Relative to Unified-IO 2, it organizes multimodal computation around recurrent thought and action conditioning rather than task unification alone. Relative to model-orchestration systems, it moves cross-modal coordination from externally programmed routing toward end-to-end learning. These are architectural differences, not evidence of achieved general intelligence; their value must be established through the implementations, ablations, and evaluation criteria developed in the following sections.
